\documentclass[11pt]{article}
\usepackage{times}
\usepackage{fullpage}
\usepackage{url}
\usepackage{hyperref}
\usepackage{fancyhdr}
\usepackage{graphicx}
\usepackage{enumitem}
\usepackage{subcaption}
\usepackage{amsmath, amsfonts, amsthm, fouriernc}
\usepackage{IEEEtrantools}
\usepackage{parskip}

\pagestyle{fancy}
\addtolength{\headheight}{2em}
\addtolength{\headsep}{1.5em}
\lhead{\large{ME 2801}}

\usepackage{sectsty}
\sectionfont{\fontsize{12}{8}\selectfont}

\begin{document}

\begin{center}
  \Large{\bf{Assignment: Frequency Response and Bode Plots}}
\end{center}

The textbook exercises designated with the word ``Nise'' are from the
``Problems'' section at the end of Chapter~10 in \emph{Control System
Engineering} by Norman Nise, 7th edition.

% ---------------------------------------------------------------
\section*{Part 1: Frequency Response by Hand}

\begin{enumerate}

\item For the transfer function $G(s) = \dfrac{10}{s + 10}$:
  \begin{enumerate}
    \item Evaluate $G(j\omega)$ at $\omega = 1$, $10$, and $100$\,rad/s.
      Compute the magnitude (in dB) and phase (in degrees) at each frequency.
    \item Sketch the straight-line (asymptotic) Bode magnitude and phase plots
      by hand.  Label the break frequency and the low- and high-frequency
      asymptote slopes.
    \item Use MATLAB \texttt{bode(G)} to plot the exact response.  Overlay your
      sketch and note where the straight-line approximation has the largest
      error.
  \end{enumerate}

\item Convert to Bode form and sketch the asymptotic Bode magnitude plot for
  each transfer function.  Label all break frequencies and asymptote slopes.
  Verify with MATLAB.

  \begin{enumerate}
    \item $G(s) = \dfrac{100(s+1)}{s(s+10)}$
    \item $G(s) = \dfrac{50}{(s+2)(s+5)}$
    \item $G(s) = \dfrac{5(s+2)}{s^2(s+20)}$
  \end{enumerate}

\end{enumerate}

% ---------------------------------------------------------------
\section*{Part 2: Bode Sketching from Building Blocks}

\begin{enumerate}[resume]

\item The transfer function of the USV surge plant is
  \[
    G_{\mathrm{surge}}(s) = \frac{5.1}{1.3\,s + 1}
  \]
  \begin{enumerate}
    \item Convert to Bode form.  Identify the DC gain (in dB) and the break
      frequency.
    \item Sketch the asymptotic Bode magnitude and phase plots.
    \item Use MATLAB to produce the exact Bode plot and compare.
    \item At what frequency does the magnitude drop to $-3$\,dB relative to DC?
      What is the physical interpretation of this frequency?
  \end{enumerate}

\end{enumerate}

% ---------------------------------------------------------------
\section*{Part 3: Nise Problems}

\begin{enumerate}[resume]

\item \textbf{Nise 10.1} (selected parts): Sketch the Bode magnitude and phase
  plots for the given transfer functions.

\item \textbf{Nise 10.2} (selected parts): From the given Bode plots, identify
  the transfer function.

\end{enumerate}

% ---------------------------------------------------------------
\section*{Part 4: MATLAB Exploration}

\begin{enumerate}[resume]

\item Use the \texttt{sketchbode} utility (in \texttt{matlab/}) to overlay the
  asymptotic and exact Bode plots for the surge plant.  Run the function and
  include the resulting figure in your report.

\item For the closed-loop surge system with $K_p = 2$, use MATLAB to plot the
  Bode magnitude of the closed-loop transfer function.  Identify the
  $-3$\,dB bandwidth.

\end{enumerate}

\end{document}
